\section{Graphical separation and adjustment: finite proofs}
\label{sec:annex-p16-graphical-separation-adjustment}

These proofs use the finite graph and structural kernels of
\cref{sec:p16-finite-causal-models}. The source's path-blocking criterion
and adjustment expression are identified in
\cref{sec:p16-graphical-separation-adjustment}; the constructions below
supply their stated finite justification.

\begin{annexitem}{annex:p16.active-walk-reduction}{thm:p16-active-walk-reduction}{Removing repetitions without closing the walk}
\begin{proof}
A path is a walk, so only the converse needs proof. Choose an active walk
from \(A\) to \(B\) of minimum length; a minimum exists among the
nonnegative integer lengths. Write its vertices as
\(i_0,\ldots,i_m\). Suppose \(i_r=i_s=w\) for \(r<s\).
Delete the portion strictly after the first occurrence through the second
occurrence, retaining the first occurrence and its new next edge from
the old second occurrence. The resulting sequence is a shorter walk
with the same endpoints. Its only possible change of internal collider
status is at the retained occurrence of \(w\).

If that occurrence is now an endpoint, there is no condition to check
there. Otherwise both former occurrences were internal. If \(w\in C\),
both were colliders, so the retained entering and leaving edges both
point into \(w\); it remains an active collider. If
\(w\in\operatorname{an}_G(C)\setminus C\), either new status is
allowed: a collider is activated and a noncollider is unconditioned.

It remains to consider \(w\notin\operatorname{an}_G(C)\).
Both former occurrences were noncolliders. The only forbidden new
status would be a collider. Suppose it occurs. The entering edge at
the first occurrence points into \(w\), so that occurrence's former
leaving edge must point away from \(w\). Likewise, the leaving edge
at the second occurrence points into \(w\), so its former entering
edge must point away from \(w\). Along the deleted portion, the
first edge is therefore oriented in the direction of travel and the
last edge is oriented against that direction. Take the first change
from an edge oriented with the travel to one oriented against it.
Their common vertex is a collider occurrence on the original walk.
The initial portion up to that vertex is a directed path from \(w\):
it cannot repeat a vertex because the graph is acyclic. Activity of
the original walk puts this collider in \(\operatorname{an}_G(C)\).
Concatenating the directed paths makes \(w\) an ancestor of \(C\),
a contradiction.

Thus the shortened walk is active in every case, contradicting
minimality. No vertex repeats on a shortest active walk, so it is an
active path. If either endpoint set is empty, neither kind of sequence
exists, and the equivalence also holds.
\end{proof}
\end{annexitem}

\begin{annexitem}{annex:p16.moral-separation-bridge}{thm:p16-moral-separation-bridge}{The exact bridge between the two graphical criteria}
\begin{proof}
Put \(H=\operatorname{an}_G(A\cup B\cup C)\). This set is
ancestral: a parent of one of its vertices has a directed path to the
same endpoint set after appending the parent's edge. The possible
zero-length path at an endpoint causes no exception.

\paragraph{An active path gives a moral connection.}
Every collider on an active path belongs to
\(\operatorname{an}_G(C)\subseteq H\). Every other internal
vertex has an outgoing edge on at least one side of the path. Follow
that edge and continue along the path as long as the arrows point in
the direction of travel. Since the path is finite, this reaches an
endpoint or a collider. The starting vertex is an ancestor of that
endpoint or collider, and hence belongs to \(H\). The path's
endpoints belong to \(H\) by definition. Thus the whole path lies
in the induced ancestral graph.

Any internal vertex on the path that belongs to \(C\) must be a
collider. Its two neighboring vertices are co-parents of that vertex,
so the moral graph joins them. Bypass every such collider by this
moral edge. Two collider occurrences cannot be adjacent on a path in
a DAG: their common edge would have to point into both ends. Thus
these bypasses leave a sequence of adjacent, distinct vertices with
no vertex in \(C\), and the original endpoints remain. It is an
undirected path in \(G_H^{\mathrm m}\setminus C\).

\paragraph{A moral connection gives an active walk.}
Conversely, take an undirected path from \(A\) to \(B\) in
\(G_H^{\mathrm m}\setminus C\). For each of its edges, retain
an original directed edge if one is present, traversing it in the
needed direction. Otherwise choose a child \(k\in H\) witnessing
the co-parent edge, and replace that edge by the two-edge sequence
\(i\to k\leftarrow j\). The resulting sequence is a walk in
the original DAG. All its vertices lie in \(H\). The vertices of
the original moral path avoid \(C\), and every inserted occurrence
is a collider. Therefore every occurrence of a vertex in \(C\)
on this walk is a collider.

This walk may contain colliders that are ancestors of \(A\) or
\(B\) but not of \(C\). Call such an occurrence bad. All other
collider occurrences are activated by \(C\), and all noncolliders
are outside \(C\). If a bad occurrence has vertex \(w\), then
\(w\in H\setminus\operatorname{an}_G(C)\), so there is a
possibly zero-length directed path from \(w\) to some \(a\in A\)
or some \(b\in B\).
Every vertex on that directed path is outside \(C\), since otherwise
\(w\) would be an ancestor of \(C\).

If the endpoint is \(a\in A\), replace the portion of the walk
from its first endpoint through this occurrence of \(w\) by the
directed path from \(w\) to \(a\), traversed backwards. Retain
the old portion from \(w\) to the \(B\) endpoint. The new internal
occurrences on the directed portion are noncolliders outside \(C\).
At its join, \(w\) has a tail on the new edge and a head on its
retained old edge, so it is now a noncollider. If the directed path
has length zero, \(w\) is the new endpoint and needs no collider
condition. All retained old internal occurrences keep their status.
The case of an endpoint \(b\in B\) is the reverse construction:
keep the old initial portion through \(w\), and follow the directed
path to \(b\).

Each replacement removes at least the chosen bad occurrence and
introduces none. It may repeat a vertex already elsewhere on the walk;
that is permitted, because collider status is local to an occurrence.
The number of bad occurrences strictly decreases, so finitely many
replacements give an active walk from \(A\) to \(B\). By
\cref{thm:p16-active-walk-reduction} an active path exists. The two
constructions prove equivalence of existence of the two connections;
negating that equivalence proves the stated separation criterion.
\end{proof}
\end{annexitem}

\begin{annexitem}{annex:p16.ancestral-marginal}{thm:p16-ancestral-marginal}{Summing away the nonancestors}
\begin{proof}
Choose a topological order of \(G\) using
\cref{thm:p16-causal-topological-order}. Sum over the coordinates
outside \(H\) in reverse topological order. When a vertex
\(j\notin H\) is reached, every child outside \(H\) has already
been removed, and it has no child in \(H\): a parent of a vertex
of an ancestral set also belongs to that set. The only remaining
factor involving \(v_j\) is consequently its own kernel
\(k_j(v_j\mid v_{\operatorname{pa}_G(j)})\).
For any fixed values of its parents, summing this factor over \(S_j\)
gives one. All other remaining factors are independent of that summation
coordinate, and finite distributivity permits them to be pulled out.
The summation therefore deletes exactly the factor indexed by \(j\).

Induction along this reverse order removes all and only the factors
outside \(H\). Every parent of a vertex in \(H\) lies in \(H\),
so the retained product is well defined on \(S_H\) and is precisely
the displayed marginal. If \(H=\varnothing\), all factors are
removed and the result is one. No division and no positive-support
assumption is used.
\end{proof}
\end{annexitem}

\begin{annexitem}{annex:p16.finite-global-markov}{thm:p16-finite-global-markov}{Separating the finite ancestral product}
\begin{proof}
An empty endpoint set gives the conditional-independence identity
immediately, so assume \(A,B\ne\varnothing\). By
\cref{thm:p16-moral-separation-bridge}, the graph
\(G_H^{\mathrm m}\setminus C\), where
\(H=\operatorname{an}_G(A\cup B\cup C)\), separates \(A\)
from \(B\). Partition its vertices into sets \(D,F\), taking
\(D\) to be the vertices reachable from \(A\) by an undirected
walk, including a walk of length zero, and \(F\) the remaining
vertices. Erasing portions between repeated vertices converts an
undirected walk to a path, so no vertex of \(B\) is in \(D\).
No edge joins \(D\) to \(F\), because appending that edge to a
walk would put its other endpoint in \(D\). Thus \(A\subseteq D\),
\(B\subseteq F\), \(H=D\cup F\cup C\) is a disjoint union,
and no moral edge joins \(D\) to \(F\).

The family \(\{i\}\cup\operatorname{pa}_G(i)\) of each
\(i\in H\) is a clique in \(G_H^{\mathrm m}\): a parent is
joined to the child by their original edge, and distinct parents are
joined through their common child \(i\in H\). Its members outside
\(C\) therefore lie entirely in \(D\) or entirely in \(F\).
Allocate its kernel to a first product if those members lie in \(D\),
and to a second product if they lie in \(F\). Allocate a kernel
whose whole family lies in \(C\) to the first product. Denote the
two nonnegative products by \(\phi(d,c)\) and \(\psi(f,c)\).
An empty product is one. The ancestral-marginal lemma gives
\[
 p_H(d,f,c)=\phi(d,c)\psi(f,c).
\]
These are functions on finite assignment sets, not newly asserted
probability laws individually.

For fixed \(a,b,c\), sum over the coordinates in \(D\setminus A\)
and \(F\setminus B\). Define
\[
 g(a,c)=\sum_{d:d_A=a}\phi(d,c),\qquad
 h(b,c)=\sum_{f:f_B=b}\psi(f,c).
\]
Finite distributivity gives
\[
 p_{A\cup B\cup C}(a,b,c)=g(a,c)h(b,c).
\]
Put \(g_C(c)=\sum_a g(a,c)\) and
\(h_C(c)=\sum_b h(b,c)\). Summing the preceding identity yields
\(p_C(c)=g_C(c)h_C(c)\). If \(p_C(c)>0\), both factors are
strictly positive, since they are nonnegative finite numbers. Dividing
the joint identity and its marginal sums by this product gives the
conditional masses, abbreviated here as \(p(a,b\mid c)\),
\(p(a\mid c)\), and \(p(b\mid c)\):
\[
\begin{aligned}
 p(a,b\mid c)&=\frac{g(a,c)}{g_C(c)}
                         \frac{h(b,c)}{h_C(c)},\\
 p(a\mid c)&=\frac{g(a,c)}{g_C(c)},\qquad
 p(b\mid c)=\frac{h(b,c)}{h_C(c)}.
\end{aligned}
\]
This proves the required conditional independence on every positive
stratum. At a null stratum no conditional expression is required.
The factorization theorem for the finite independent-disturbance model
supplies exactly the kernels assumed here, proving the final assertion.
\end{proof}
\end{annexitem}

\begin{annexitem}{annex:p16.finite-back-door-adjustment}{thm:p16-finite-back-door-adjustment}{A modified model proves the adjustment identity}
\begin{proof}
Fix the value \(x\in S_t\). Write \(G_{\underline t}\) for the
DAG obtained by deleting every edge whose initial vertex is \(t\),
so all outgoing edges from the treatment are removed. This is not
the intervention graph, which deletes incoming edges. Define an
auxiliary structural model \(M_x^\dagger\) on the same probability
space, disturbances, and endogenous value sets as follows.
Keep the original equation at \(t\). For each original child of
\(t\), replace the argument supplied by that parent with the fixed
value \(x\), and remove that parent from the function's list of
variable inputs. Keep all other functions unchanged. The resulting
functions are total, the allowed graph is \(G_{\underline t}\),
and the disturbances remain mutually independent. Denote its solution
by \(W\).

\paragraph{Pointwise comparison of the three solutions.}
Let \(R=N\setminus(\{t\}\cup\operatorname{de}_G(\{t\}))\).
Every parent of a vertex in \(R\) is again in \(R\), and no
equation there is altered in either modified model. Topological
evaluation therefore gives
\(W_R=V_R=V_R^{t=x}\) at every outcome, as in the proof of
\cref{thm:p16-causal-intervention-laws}. Every parent of \(t\)
belongs to \(R\), because a directed path from \(t\) to one of
its parents would create a cycle. The retained equation at \(t\)
thus gives \(W_t=V_t=X\) pointwise.

For every vertex other than \(t\), compare the auxiliary solution
with the intervention solution in a topological order. If the vertex
was a child of \(t\), both of its evaluations use the prescribed
value \(x\) in that parent position. Every other parent value agrees
by induction. Its disturbance is unchanged, so its output agrees.
At a vertex with no parent \(t\), all parent values and its function
agree directly. Consequently
\[
 W_t=X,\qquad W_B=Y^{t=x},\qquad W_C=Z^{t=x}=Z.
\]
The last equality also uses the condition that \(C\) consists of
nondescendants and excludes \(t\). These identities include outcomes
of zero probability.

\paragraph{The graphical implication in the auxiliary model.}
Any path from \(t\) to \(B\) in \(G_{\underline t}\) starts
with an arrow into \(t\), and is a back-door path in the original
graph. Its edge directions, colliders, and noncolliders are unchanged.
If it were active given \(C\) in \(G_{\underline t}\), every
collider would have a directed path to \(C\) in that graph and
hence also in \(G\). It would then be an active back-door path
in \(G\), contradicting the hypothesis. Thus
\(\{t\}\mathrel{\perp_{G_{\underline t}}}B\mid C\).
Apply \cref{thm:p16-finite-global-markov} to the auxiliary model's
finite factorization. The pointwise identities give
\[
 X\mathrel{\perp\!\!\!\perp}Y^{t=x}\mid Z.
\]
No transformation rule for do expressions has been assumed.

\paragraph{Positive conditioning and consistency.}
Take \(z\) with \(\mathbb P(Z=z)>0\). Overlap gives
\(\mathbb P(X=x,Z=z)>0\), so all of the following ratios exist.
Conditional independence, divided by
\(\mathbb P(X=x\mid Z=z)>0\), gives
\[
 \mathbb P(Y^{t=x}=y\mid Z=z)
   =\mathbb P(Y^{t=x}=y\mid X=x,Z=z).
\]
On the event \(X=x\), pointwise consistency in
\cref{thm:p16-causal-intervention-laws} states
\(V^{t=x}=V\). In particular the events in the numerators on
that event agree, giving
\[
 \mathbb P(Y^{t=x}=y\mid X=x,Z=z)
    =\mathbb P(Y=y\mid X=x,Z=z).
\]
Since \(Z^{t=x}=Z\) pointwise, the left side of the theorem's
stratum assertion is the first probability above. Partitioning the
probability space by the finitely many values of \(Z\) now gives
\[
 \mathbb P(Y^{t=x}=y)
 =\sum_{z:\mathbb P(Z=z)>0}
       \mathbb P(Y=y\mid X=x,Z=z)\mathbb P(Z=z).
\]
Null strata have no mass in this partition and require no conditional
value. For empty \(C\), the same argument has one stratum of mass
one, and overlap is simply \(\mathbb P(X=x)>0\).

Finally, each factor on the right is a marginal sum or a ratio of
marginal sums of \(p_{M,O}\), because all coordinates in
\(\{t\}\cup B\cup C\) are observed. For any two models in the
stated class with the same observed law, the positive strata, their
weights, and all those ratios agree. Their intervention laws therefore
agree at every \(y\). This is identification according to
\cref{def:p16-causal-identification}, even if coordinates outside
\(O\) are unobserved. It establishes no finite-sample estimation
claim.
\end{proof}
\end{annexitem}

\begin{annexitem}{annex:p16.adjustment-numerical-example}{thm:p16-adjustment-numerical-example}{Every mass and both adjusted intervention values}
\begin{proof}
The exogenous weights sum to
\[
 \left(\tfrac12+\tfrac12\right)
       \left(\tfrac34+\tfrac14\right)
       \left(\sum_{u=1}^8\tfrac18\right)=1
\]
and factor over the three coordinates, proving their mutual
independence. In topological order \(Z,X,Y\), the equations are
total and evaluate to the displayed assignments. Their kernels are
\[
 k_Z(z)=\tfrac12,\quad k_X(x\mid z)=\alpha(x\oplus z),\quad
 k_Y(1\mid x,z)=\frac{1+2x+4z}{8}.
\]
The four thresholds are \(1,3,5,7\), so the final equality counts
exactly that many of the eight equally likely disturbance values.
The zero-output kernel is one minus this value.

For \((z,x)=(0,0)\), the first two factors multiply to \(3/8\),
and the response probabilities are \(7/8,1/8\); their products
are \(21/64,3/64\). For \((0,1)\), the parent mass is \(1/8\)
and the response probabilities are \(5/8,3/8\), giving
\(5/64,3/64\). For \((1,0)\), the parent mass is \(1/8\)
and the response probabilities are \(3/8,5/8\), giving
\(3/64,5/64\). Finally, for \((1,1)\), the parent mass is
\(3/8\) and the response probabilities are \(1/8,7/8\), giving
\(3/64,21/64\). This proves all eight masses in the body table;
they sum to \(64/64\).

Summing over \(Z\), the \((X,Y)\) table is
\[
\begin{array}{c|cc|c}
 &Y=0&Y=1&\text{row total}\\\hline
 X=0&3/8&1/8&1/2\\
 X=1&1/8&3/8&1/2
\end{array}
\]
Division by the positive row masses gives the two conditional response
probabilities \(1/4\) and \(3/4\), and their difference \(1/2\).

Under equation replacement \(X=x\), the unchanged fair \(Z\)
and uniform \(U_Y\) give, directly,
\[
 \mathbb P(Y^{X=x}=1)
 =\tfrac12\frac{1+2x}{8}
       +\tfrac12\frac{5+2x}{8}
 =\frac{3+2x}{8}.
\]
Here a superscript \(X=x\) names intervention at the vertex whose
variable is called \(X\). The two values are \(3/8\) and
\(5/8\), whose difference is \(1/4\).

The only back-door path in this three-vertex graph is
\(X\leftarrow Z\to Y\). Conditioning on its noncollider \(Z\)
blocks it, and \(Z\) is not a descendant of \(X\).
For every \(x,z\),
\(\mathbb P(X=x,Z=z)=\frac12\alpha(x\oplus z)>0\).
The conditional response ratio is therefore
\((1+2x+4z)/8\). Its weighted sum with \(\mathbb P(Z=z)=1/2\)
is exactly the preceding intervention calculation. This checks both
the graph and every positive denominator used for adjustment.
\end{proof}
\end{annexitem}

\begin{annexitem}{annex:p16.collider-numerical-example}{thm:p16-collider-numerical-example}{Selection, intervention, and an active-path counterexample}
\begin{proof}
The four equally likely disturbance pairs evaluate to the following
four full assignments, each of mass \(1/4\):
\[
\begin{array}{cccc|c}
 X&Y&K&D&\text{mass}\\\hline
 0&0&0&0&1/4\\
 0&1&1&1&1/4\\
 1&0&1&1&1/4\\
 1&1&1&1&1/4
\end{array}
\]
Thus the marginal \((X,Y)\) table is uniform, and each variable
has probabilities \((1/2,1/2)\); the joint masses are the products
of their marginal masses. The event \(K=1\) has mass \(3/4\).
Conditioning divides each of its three assignment masses by \(3/4\),
giving \(1/3\) at the three pairs other than \((0,0)\).
Under this conditional law,
\[
 \mathbb P(X=1\mid K=1)=\mathbb P(Y=1\mid K=1)=2/3,
 \qquad \mathbb P(X=1,Y=1\mid K=1)=1/3\ne4/9.
\]
This is dependence on a positive stratum and disproves conditional
independence given \(K\). On the other positive stratum \(K=0\),
both variables are constantly zero; there they are independent.
The definition requires the factorization on every positive stratum,
so this does not restore conditional independence given \(K\).
Since \(D=K\) pointwise, conditioning on \(D=1\) gives exactly
the same table and dependence.

The path \(X\to K\leftarrow Y\) has its sole internal vertex
as a collider. With empty conditioning set it is blocked, because
\(K\) has no descendant in that empty set. Given \(\{K\}\)
the collider itself is conditioned on; given \(\{D\}\) its
descendant is conditioned on. In each case it is active. There is
no other path from \(X\) to \(Y\) in this graph.

Intervening on \(K\) replaces its equation by \(K=1\), and the
unchanged equation for \(D\) then gives \(D=1\). Neither the
equations for \(X,Y\) nor their outcome weights change, so all
four \((X,Y)\) pairs still have mass \(1/4\). This establishes
the asserted contrast with conditioning.

For the final counterexample, take three independent fair bits
\(U_X,U_Y,U_K\), a constant disturbance for \(D\), and the
same allowed graph and binary endogenous value sets. Put
\[
 X=U_X,\qquad Y=U_Y,\qquad K=U_K,\qquad D=K.
\]
The function at \(K\) ignores its allowed parent arguments,
which \cref{def:p16-finite-scm} explicitly permits. The eight triples
\((X,Y,K)\) have mass \(1/8\), while \(D\) is their last
coordinate. For either \(k=0\) or \(k=1\), the event \(K=k\)
has mass \(1/2\), and its conditional \((X,Y)\) masses are
\((1/8)/(1/2)=1/4\). The marginals are both \((1/2,1/2)\),
so conditional independence holds on both positive strata. The allowed
graph is unchanged and its path is still active given \(\{K\}\).
It follows that nonseparation alone cannot force dependence in every
compatible model. No faithfulness assumption was used or supplied.
\end{proof}
\end{annexitem}
